%% ============================================================
%%  Defensive Publication — Prior Art Disclosure
%%  PIRDS on MOC: Prime-Indexed Resonant Dynamical Systems
%%  as a Governed Bridge from Multiplicity Theory to
%%  Trainable Spectral Learning
%%
%%  Author:  Multiplicity Theorist / PhaseMirror Project
%%  Date:    April 2026
%%  Status:  Defensive Publication / Prior Art Disclosure
%% ============================================================

\documentclass[11pt,a4paper]{article}

%% ---- Packages -----------------------------------------------
\usepackage[T1]{fontenc}
\usepackage[utf8]{inputenc}
\usepackage{lmodern}
\usepackage{microtype}
\usepackage{geometry}
\geometry{
  a4paper,
  top=28mm, bottom=32mm,
  left=30mm, right=30mm,
  headheight=14pt
}

\usepackage{amsmath,amssymb,amsthm}
\usepackage{mathtools}
\usepackage{bm}
\usepackage{physics}
\usepackage{xcolor}
\usepackage{booktabs}
\usepackage{tabularx}
\usepackage{array}
\usepackage{multirow}
\usepackage{enumitem}
\usepackage{graphicx}
\usepackage{float}
\usepackage{fancyhdr}
\usepackage{titlesec}
\usepackage{tocloft}
\usepackage{hyperref}
\usepackage{cleveref}
\usepackage{listings}
\usepackage{mdframed}
\usepackage{tcolorbox}
% fontawesome5 not available in this environment
\tcbuselibrary{breakable,skins}
\usepackage{caption}
\usepackage{subcaption}
\usepackage{algorithmicx}
\usepackage{algpseudocode}
\usepackage[ruled,vlined]{algorithm2e}

%% ---- Colour palette -----------------------------------------
\definecolor{teal}{HTML}{01696F}
\definecolor{tealdark}{HTML}{0C4E54}
\definecolor{stone}{HTML}{28251D}
\definecolor{muted}{HTML}{7A7974}
\definecolor{border}{HTML}{D4D1CA}
\definecolor{bg}{HTML}{F7F6F2}
\definecolor{rust}{HTML}{964219}
\definecolor{codebase}{HTML}{1C1B19}
\definecolor{codekw}{HTML}{4F98A3}
\definecolor{codestr}{HTML}{6DAA45}
\definecolor{codecomm}{HTML}{797876}

%% ---- Listings (Python code) ---------------------------------
\lstset{
  language=Python,
  basicstyle=\ttfamily\small,
  keywordstyle=\color{codekw}\bfseries,
  stringstyle=\color{codestr},
  commentstyle=\color{codecomm}\itshape,
  backgroundcolor=\color{bg},
  frame=single,
  framerule=0.4pt,
  rulecolor=\color{border},
  numbers=left,
  numberstyle=\tiny\color{muted},
  numbersep=6pt,
  breaklines=true,
  captionpos=b,
  showstringspaces=false,
  tabsize=4,
  xleftmargin=2em,
}

%% ---- Theorem environments -----------------------------------
\theoremstyle{plain}
\newtheorem{theorem}{Theorem}[section]
\newtheorem{proposition}[theorem]{Proposition}
\newtheorem{lemma}[theorem]{Lemma}
\newtheorem{corollary}[theorem]{Corollary}
\theoremstyle{definition}
\newtheorem{definition}[theorem]{Definition}
\newtheorem{example}[theorem]{Example}
\theoremstyle{remark}
\newtheorem{remark}[theorem]{Remark}
\newtheorem{claim}[theorem]{Claim}

%% ---- Section styling ----------------------------------------
\titleformat{\section}{\large\bfseries\color{teal}}{{\thesection}}{0.75em}{}[\color{border}\hrule\vspace{2pt}]
\titleformat{\subsection}{\normalsize\bfseries\color{stone}}{{\thesubsection}}{0.6em}{}
\titleformat{\subsubsection}{\normalsize\itshape\color{stone}}{{\thesubsubsection}}{0.5em}{}

%% ---- Header / footer ----------------------------------------
\pagestyle{fancy}
\fancyhf{}
\fancyhead[L]{\small\color{muted}\textit{PIRDS on MOC — Defensive Publication}}
\fancyhead[R]{\small\color{muted}April 2026}
\fancyfoot[C]{\small\color{muted}\thepage}
\renewcommand{\headrulewidth}{0.3pt}
\renewcommand{\headrule}{\color{border}\hrule height\headrulewidth\hfill}

%% ---- Hyperref setup -----------------------------------------
\hypersetup{
  colorlinks=true,
  linkcolor=teal,
  citecolor=teal,
  urlcolor=teal,
  pdftitle={PIRDS on MOC: Defensive Publication},
  pdfauthor={PhaseMirror Project},
  pdfsubject={Prime-Indexed Resonant Dynamical Systems},
}

%% ---- Custom commands ----------------------------------------
\newcommand{\primes}{\mathcal{P}}
\newcommand{\tiers}{\mathcal{D}}
\newcommand{\multsp}{\mathcal{M}}
\newcommand{\Sp}{S_p}
\newcommand{\Apr}{A_{p^r}^{\alpha}}
\newcommand{\Rpr}{R_{p^r}^{\varphi}}
\newcommand{\Pip}{\Pi_{p^r}}
\newcommand{\Delp}{\Delta_{p^r}}
\newcommand{\Xn}{X_n}
\newcommand{\Tn}{T_n}
\newcommand{\Up}{U_p}
\newcommand{\Orth}{\Omega_{\mathrm{orth}}}
\newcommand{\Stab}{\Omega_{\mathrm{stab}}}
\newcommand{\ResF}{\mathcal{R}}
\newcommand{\state}{M}
\newcommand{\xhat}{\hat{x}}
\newcommand{\Kd}{K_d}
\newcommand{\Kdp}{K'_d}
\newcommand{\Loss}{\mathcal{L}}
\newcommand{\faithful}{\textcolor{teal}{\textbf{FAITHFUL}}}
\newcommand{\extension}{\textcolor{rust}{\textbf{EXTENSION}}}
\newcommand{\corrected}{\textcolor{muted}{\textbf{CORRECTED}}}

%% ---- tcolorbox styles --------------------------------------
\tcbset{
  defbox/.style={
    colback=bg, colframe=teal,
    fonttitle=\bfseries\small\color{white},
    coltitle=white,
    attach boxed title to top left={yshift=-2mm,xshift=4mm},
    boxed title style={colback=teal, colframe=teal},
    breakable, enhanced,
    top=4mm, bottom=2mm, left=3mm, right=3mm,
  },
  warnbox/.style={
    colback=bg, colframe=rust,
    fonttitle=\bfseries\small\color{white},
    coltitle=white,
    attach boxed title to top left={yshift=-2mm,xshift=4mm},
    boxed title style={colback=rust, colframe=rust},
    breakable, enhanced,
    top=4mm, bottom=2mm, left=3mm, right=3mm,
  },
  notebox/.style={
    colback=bg, colframe=border,
    breakable, enhanced,
    top=3mm, bottom=3mm, left=3mm, right=3mm,
  }
}

%% ============================================================
\begin{document}
%% ============================================================

%% ---- Title page --------------------------------------------
\begin{titlepage}
\pagecolor{bg}
\vspace*{2cm}
\begin{center}
  {\color{teal}\rule{0.9\textwidth}{2pt}}\\[1.2em]
  {\fontsize{22}{28}\selectfont\bfseries\color{stone}
    Prime-Indexed Resonant Dynamical Systems\\[0.4em]
    as a Governed Bridge from\\[0.4em]
    Multiplicity Operator Calculus\\[0.4em]
    to Trainable Spectral Learning}\\[1.0em]
  {\color{teal}\rule{0.9\textwidth}{2pt}}\\[2.0em]
  {\large\bfseries\color{muted} DEFENSIVE PUBLICATION --- PRIOR ART DISCLOSURE}\\[0.5em]
  {\normalsize\color{muted}
    \textit{All technical content herein is disclosed for the purpose of establishing prior art.}\\
    \textit{No patent rights are asserted. This document enters the public domain upon release.}}
  \\[2.5em]
  {\large\bfseries\color{stone} PhaseMirror Project}\\[0.3em]
  {\normalsize\color{muted} Multiplicity Theory Research Group}\\[0.3em]
  {\normalsize\color{muted} April 2026}\\[3.5em]
  \begin{tabular}{ll}
    \textbf{Repository:} & \texttt{PhaseMirror/PhaseMirror-HQ}\\[0.3em]
    \textbf{Core modules:} & \texttt{moc\_resonance.py},\ \texttt{pirds\_module.py},\ \texttt{moc\_108\_repro.py}\\[0.3em]
    \textbf{ADRs:} & \texttt{ADR-001},\ \texttt{ADR-002-pirds-maps-to-moc}\\[0.3em]
    \textbf{Status:} & Public prior art disclosure\\[0.3em]
    \textbf{Version:} & 1.0\\
  \end{tabular}
  \\[3.5em]
  {\color{teal}\rule{0.9\textwidth}{0.8pt}}\\[0.8em]
  {\small\color{muted}
    \textit{This document records the architecture, mathematics, and implementation of the\\
    PIRDS framework as developed in the PhaseMirror research programme during\\
    April 2026. The purpose is to permanently and publicly establish prior art,\\
    prevent subsequent monopolization of these methods, and invite open scientific scrutiny.}}
\end{center}
\end{titlepage}
\nopagecolor

%% ---- Executive Summary --------------------------------------
\newpage
\begin{tcolorbox}[defbox, title={Executive Summary}]
\noindent\textbf{What is disclosed.}
This document discloses the complete design of a \emph{Prime-Indexed Resonant Dynamical System} (PIRDS), a mathematical framework and executable runtime that unifies \emph{Multiplicity Operator Calculus} (MOC) with gradient-based spectral learning. The system treats any cyclically-structured signal space as a state space governed by prime-indexed operator families, and measures structural fidelity through a three-component resonance functional derived from CRT tier analysis.

\vspace{0.5em}
\noindent\textbf{Key innovations disclosed.}
\begin{enumerate}[leftmargin=1.5em,itemsep=2pt]
  \item A faithful executable implementation of the MOC operator algebra: subdivision $S_p$, accent $A_{p^r}^\alpha$, rotation $R_{p^r}^\varphi$, and CRT projectors $\Pi_{p^r}$, $\Delta_{p^r}$.
  \item The exact MOC resonance functional $R = \lambda_1 R_1 + \lambda_2 R_2 + \lambda_3 R_3$ with Hann-tapered shift-optimised $R_1$, comb-index geometric-mean $R_2$, and circular phase-coherence $R_3$.
  \item A trainable spectral basis $U(\omega_p) = Q\,\mathrm{diag}(\varphi(\omega_p))\,Q^{-1}$ with learnable prime weights $w_p(t)$ and orthogonality regulariser $\Omega_{\mathrm{orth}}$.
  \item The C4 prime-grid ablation protocol: a pre-registered, falsifiable experimental design that tests whether prime-indexed structure is structurally necessary versus aesthetically motivated.
  \item The Q-basis ablation study: a governance mechanism bounding when a learned (Koopman-style) basis $Q$ may replace the fixed DFT.
  \item Complete identification of PhaseMirror-HQ as the first concrete PIRDS instantiation, with the MOC document as its canonical mathematical kernel.
\end{enumerate}

\vspace{0.5em}
\noindent\textbf{Prior art claim.}
All methods, architectures, protocols, and code patterns described herein were conceived and implemented by the PhaseMirror Project during April 2026. This disclosure establishes that prior art, prevents subsequent patent filings on these specific combinations of prime-indexed operator calculus and gradient-based spectral learning, and places the work in the permanent public record.
\end{tcolorbox}

\tableofcontents
\newpage

%% =============================================================
\section{Introduction and Motivation}
%% =============================================================

\subsection{Context: PhaseMirror-HQ as a PIRDS}

The PhaseMirror-HQ repository is not merely an implementation that happens to use
prime-indexed structure --- it \emph{is} a Prime-Indexed Resonant Dynamical System.
The repository's evolution is therefore not the addition of a new component, but the
generalisation of an existing PIRDS into a foundational, reusable operating framework.

This recognition crystallised through the identification that the Multiplicity Operator
Calculus (MOC) document --- previously treated as a theoretical companion --- already
constitutes a complete Prime-Indexed Operator Dynamics layer. Specifically, MOC provides:
prime-indexed operator families $S_p, A_{p^r}^\alpha, R_{p^r}^\varphi, \Pi_{p^r},
\Delta_{p^r}$ acting on a multiplicity state $M = (X, H)$; a three-component resonance
functional $\Res \in [0,1]$ as the structural fidelity measure; conservation invariants
enforced via projection; and a noncommutative composite-word evolution object
$\bar{N} = \prod_{p \mid n} \hat{P}^{(p)}$ as the discrete-time dynamics.

The conclusion drawn --- and disclosed here as prior art --- is that \textbf{PIRDS $=$ MOC core $+$ learnable delta}, where the delta consists of: learnable prime weights $w_p(t)$, an explicit shared spectral basis $Q$ with trainable $\varphi(\omega_p)$, and an orthogonality regulariser $\Omega_\mathrm{orth}$. A fourth addition is not mathematical but operational: the C4 prime-grid ablation protocol, which transforms the structural claims of MOC from assertion to falsifiable experiment.

\subsection{Purpose of This Document}

This report serves three simultaneous roles:

\begin{enumerate}[leftmargin=1.5em,itemsep=3pt]
  \item \textbf{Defensive publication.} A permanent, dated, public record of every technical
    idea, architecture decision, mathematical formulation, and code pattern developed in the
    PIRDS/MOC programme during April 2026. This blocks subsequent patent applications on these
    combinations and places the ideas irrevocably in the commons.
  \item \textbf{Technical specification.} A comprehensive reference for every mathematical
    object, its relationship to prior MOC definitions, its faithful or extension status, and
    its governance constraints.
  \item \textbf{Research programme roadmap.} The three bounded studies (108-cycle
    reproduction, C4 ablation, Q-basis ablation) and the ADR governance documents establish
    the ordered sequence of experiments that will elevate PIRDS from a prototype to a
    validated, falsifiable system.
\end{enumerate}

\subsection{Central Tension}

The central tension of the programme, now precisely stated, is:

\begin{tcolorbox}[notebox]
  \textit{Can an implementation remain attached to MOC's structural claims --- specifically
  that prime-indexed tier structure is doing explanatory work --- without quietly collapsing
  into generic spectral machine learning?}
\end{tcolorbox}

This tension admits a tractable resolution: maintain an explicit, continuously-updated
classification of every code component as either \faithful{} (direct transliteration of a
MOC definition) or \extension{} (new addition that may be removed without invalidating MOC);
and run the C4 ablation before asserting structural necessity of primes.

%% =============================================================
\section{Mathematical Foundations: Multiplicity Operator Calculus}
%% =============================================================

\subsection{Multiplicity Spaces and State Representation}

\begin{definition}[Cycle group and signal space]
  Let $n \in \mathbb{Z}_{>0}$. The \emph{cycle group} is $T_n = \mathbb{Z}/n\mathbb{Z}$.
  The \emph{signal space} is $X_n = \{x : T_n \to \mathbb{R}^k\}$, identified with
  $\mathbb{R}^{n \times k}$.
\end{definition}

\begin{definition}[Multiplicity state]
  A \emph{multiplicity state} is a pair $M = (X, H)$ where:
  \begin{itemize}[itemsep=2pt]
    \item $X \in X_n$ is the signal carrier --- a function on the cycle group;
    \item $H : E \to \mathbb{R}^m$ assigns state vectors to hyperedges $e \in E$
      of an associated hypergraph encoding relational structure.
  \end{itemize}
  The space of all multiplicity states is denoted $\multsp$.
\end{definition}

In the implementation, $M$ is represented as the named tuple
\texttt{MultipliciityState(X, H\_flat)} where \texttt{X} has shape $(B, n, k)$
and \texttt{H\_flat} has shape $(B, E, m)$, with $B$ the batch dimension.

\subsection{CRT Tier Structure}

\begin{definition}[CRT tiers]
  Let $n = \prod_{p \mid n} p^{r_p}$ be the prime factorisation of $n$.
  The \emph{CRT tier set} is
  \[
    \tiers(n) = \{ p^j : p \mid n,\ 1 \le j \le r_p \}.
  \]
  For $n = 108 = 2^2 \cdot 3^3$, we have $\tiers(108) = \{2, 4, 3, 9, 27\}$.
\end{definition}

\begin{definition}[Comb indices]
  For $d \in \tiers(n)$ with $d \mid n$, the \emph{DFT comb index set} at tier $d$ is
  \[
    K_d = \bigl\{ \ell \cdot \tfrac{n}{d} : \ell = 0, 1, \ldots, d-1 \bigr\} \subset \mathbb{Z}/n\mathbb{Z}.
  \]
  The \emph{DC-excluded comb} is $K'_d = K_d \setminus \{0\}$.
\end{definition}

The significance of $K_d$ is that frequencies in $K_d$ correspond precisely to
$d$-periodic components of $x$ --- the spectral signature of tier-$d$ periodicity.

\subsection{Prime-Indexed Operator Families}

The MOC defines five families of operators on $X_n$:

\begin{definition}[Subdivision $S_p$]
  \[
    (S_p\, x)(pt + u) = x(t), \quad t \in T_n,\ u \in \{0,\ldots,p-1\},
  \]
  mapping $X_n \to X_{pn}$. Each sample is repeated $p$ times.
\end{definition}

\begin{definition}[Accent $A_{p^r}^\alpha$]
  \[
    (A_{p^r}^\alpha\, x)(t) = x(t) + \alpha\, [p^r \mid t]\, e_1,
  \]
  where $[p^r \mid t]$ is the divisibility indicator, $e_1$ is the first basis vector
  in $\mathbb{R}^k$, and $\alpha \in \mathbb{R}$ is the accent strength.
\end{definition}

\begin{definition}[Rotation $R_{p^r}^\varphi$]
  \[
    (R_{p^r}^\varphi\, x)(t) = x\!\left(t + \varphi \cdot \tfrac{n}{p^r} \bmod n\right),
    \quad \varphi \in \mathbb{Z},
  \]
  a cyclic shift by $\varphi$ sub-period units at tier $p^r$.
\end{definition}

\begin{definition}[Averaging projector $\Pi_{p^r}$]
  \[
    (\Pi_{p^r}\, x)(t) = \frac{1}{p^r} \sum_{u=0}^{p^r - 1} x\!\left(t + u \cdot \tfrac{n}{p^r}\right).
  \]
  This is the conditional expectation onto the sub-sigma-algebra generated by the
  coset partition of $T_n$ by $\frac{n}{p^r}\mathbb{Z}/n\mathbb{Z}$.
\end{definition}

\begin{definition}[Spike gate $\Delta_{p^r}$]
  \[
    (\Delta_{p^r}\, x)(t) = [p^r \mid t]\, x(t).
  \]
  Selects only the ticks at which $p^r$ divides $t$, zeroing all others.
\end{definition}

\begin{definition}[Operator words and composite evolution]
  An \emph{operator word} of length $m$ is a sequence $W = \hat{O}_m \cdots \hat{O}_1$
  of operators from the families above, applied right-to-left. The
  \emph{composite evolution operator} at cycle $n$ is
  \[
    \bar{N} = \prod_{p \mid n} \hat{P}^{(p)},
  \]
  where $\hat{P}^{(p)}$ encodes the prime-$p$ contribution. Applying $\bar{N}$ to
  $M_t$ gives $M_{t+1}$, yielding a discrete-time dynamical system on $\multsp$.
\end{definition}

\begin{remark}[Noncommutativity]
  The operator families do not commute in general:
  $A_{p^r}^\alpha \circ S_p \neq S_p \circ A_{p^r}^\alpha$.
  Word order carries structural information --- the $W_\triangle$ (ternary-first) and
  $W_\square$ (binary-first) words in the 108-cycle example produce measurably different
  resonance profiles (see \cref{sec:108cycle}).
\end{remark}

\subsection{Gauge Equivalence and Invariants}

\begin{definition}[Gauge automorphism]
  Two operator words $W$ and $W'$ are \emph{gauge-equivalent} if $R(W\,M, D) = R(W'\,M, D)$
  for all states $M$ and all datasets $D$. The \emph{gauge group} $\mathrm{Aut}(\multsp)$
  consists of all such equivalences.
\end{definition}

Three invariants are preserved by admissible compositions and enforced via projection:
\begin{align}
  E[x] &= \sum_{t \in T_n} \|x(t)\|^2 \quad \text{(energy)}, \\
  A[\sigma, \tau] &= \sum_{e \in E} \sigma(e) \cdot \tau(e) \quad \text{(area)}, \\
  F[\sigma] &= -\sum_{e \in E} \sigma(e) \log \sigma(e) \quad \text{(fairness/entropy)}.
\end{align}

%% =============================================================
\section{The Resonance Functional}
%% =============================================================

The resonance functional $\Res(x, D) \in [0,1]$ is the core measurement object of MOC.
It quantifies how structurally similar a model output $x$ is to a dataset $D$, with
sensitivity to three distinct aspects of similarity.

\subsection{Component Definitions}

\begin{definition}[$R_1$: Temporal Cross-Correlation]
  Let $W = \mathrm{diag}(w[t])$ be the Hann taper, $s_\kappa(a) = \kappa\tanh(a/\kappa)$
  the soft-clip function ($\kappa = 0$ means identity). Then
  \[
    R_1(x, D) = \max_{\tau \in T_n}
    \frac{\langle Wx,\, W T_\tau D \rangle_2^2}
         {\|Wx\|_2^2\, \|WT_\tau D\|_2^2},
  \]
  where $T_\tau$ denotes cyclic shift by $\tau$. The squaring enforces
  non-negativity and stabilises near-ties. The maximum over $\tau$ is computed
  efficiently via FFT-based cross-correlation in $O(n \log n)$.
\end{definition}

\begin{definition}[$R_2$: Tiered Spectral Energy Matching]
  Let $\hat{x}$ denote the DFT of $x$, and let
  $E_x(K) = \sum_{k \in K} |\hat{x}[k]|^2 / \sum_{k=0}^{n-1} |\hat{x}[k]|^2$
  be the energy fraction at index set $K$. With tier weights $\eta_d \ge 0$,
  $\sum_{d \in \tiers} \eta_d = 1$:
  \[
    R_2(x, D) = \sum_{d \in \tiers(n)} \eta_d \sqrt{E_x(K_d) \cdot E_D(K_d)} \in [0,1].
  \]
  This is the geometric mean of energy fractions at CRT comb positions, not
  a time-domain proxy. It is sensitive specifically to whether $x$ and $D$
  concentrate spectral energy at the same prime-indexed frequencies.
\end{definition}

\begin{definition}[$R_3$: Phase Coherence on CRT Combs]
  \[
    R_3(x, D) = \sum_{d \in \tiers(n)} \eta_d \cdot
    \left\lvert \frac{1}{|K'_d|} \sum_{k \in K'_d}
    \exp\!\bigl(i(\arg \hat{x}[k] - \arg \hat{D}[k])\bigr) \right\rvert \in [0,1].
  \]
  The inner quantity is the \emph{circular mean modulus} of phase differences at the
  DC-excluded comb positions. A value of 1 indicates perfect phase alignment at all
  tier-$d$ comb frequencies; 0 indicates uniform phase scatter.
\end{definition}

\begin{definition}[Aggregate resonance]
  \[
    \boxed{R(x, D) = \lambda_1 R_1(x, D) + \lambda_2 R_2(x, D) + \lambda_3 R_3(x, D)}
  \]
  with default weights $(\lambda_1, \lambda_2, \lambda_3) = (0.40, 0.35, 0.25)$.
  The default tier weights are uniform: $\eta_d = 1/|\tiers(n)|$ for all $d \in \tiers(n)$.
\end{definition}

\begin{remark}[Correction of v0 approximations]
  The initial implementation (v0) used three approximations that have since been
  replaced by the exact definitions above:
  \begin{center}
  \begin{tabular}{lll}
    \toprule
    Component & v0 approximation & v1 (exact) \\
    \midrule
    $R_1$ & Plain dot product, no shift search & FFT cross-corr., max over $\tau$, Hann-tapered \\
    $R_2$ & Spike-gate $\Delta_{p^r}$ (time domain) & Comb sets $K_d$ in DFT domain, geom.~mean \\
    $R_3$ & Global cosine similarity of phase angles & Circular mean per CRT comb tier, DC excluded \\
    \bottomrule
  \end{tabular}
  \end{center}
  The $R_2$ correction is structurally significant: $\Delta_{p^r}$ selects time-domain
  ticks at multiples of $p^r$, whereas $K_d$ selects DFT frequency bins. These coincide
  only in degenerate cases and encode different information about the signal's prime structure.
\end{remark}

%% =============================================================
\section{PIRDS: The Trainable Extension}
%% =============================================================

\subsection{Core Identification}

\begin{definition}[PIRDS identification]
  A \emph{Prime-Indexed Resonant Dynamical System} (PIRDS) is identified as:
  \[
    \mathrm{PIRDS} = \mathrm{MOC\ core} + \Delta_{\mathrm{learn}},
  \]
  where $\Delta_{\mathrm{learn}}$ consists of:
  \begin{enumerate}[itemsep=2pt]
    \item Learnable prime weights $w_p(t)$ replacing hand-set accent coefficients;
    \item An explicit shared spectral basis $Q$ with trainable spectral multipliers $\varphi(\omega_p)$;
    \item An orthogonality regulariser $\Omega_\mathrm{orth}$ addressing the identifiability gap.
  \end{enumerate}
  PhaseMirror-HQ is the first concrete PIRDS instantiation.
\end{definition}

\subsection{Spectral Basis and Prime Operators}

\begin{definition}[Spectral prime operator]
  For each prime $p$, the \emph{spectral prime operator} is
  \[
    U(\omega_p) = Q\, \mathrm{diag}(\varphi(\omega_p))\, Q^{-1},
  \]
  where $Q \in \mathbb{C}^{n \times n}$ is the shared spectral basis (default: DFT),
  and $\varphi(\omega_p) \in \mathbb{C}^n$ is a complex diagonal multiplier whose
  support is determined by the CRT tiers: $\mathrm{supp}(\varphi(\omega_p)) \subseteq K_p$.
  The spectral positions are indexed as $\omega_{p^r} \propto \log(p^r)$.
\end{definition}

The default $Q$ is the DFT matrix, making $U(\omega_p)$ a frequency-domain filter
whose pass-band is confined to the prime-$p$ comb $K_p$. The option \texttt{learnable\_Q=True}
allows $Q$ to be trained jointly with $\varphi$, yielding a Koopman-style learned basis ---
but this mode is subject to strict governance constraints (see \cref{sec:governance}).

\subsection{Learnable Prime Weights}

\begin{definition}[Prime weight network]
  A \emph{prime weight network} $w_\theta : \mathbb{R}^c \to \Delta^{P-1}$ is a
  parametric function (static logits or shallow MLP) mapping context embeddings to
  a probability simplex over $P$ primes:
  \[
    w_p(t) = \mathrm{softmax}(\theta_p \text{ or } \mathrm{MLP}_\theta(\text{ctx}(t)))_p,
    \quad \sum_p w_p(t) = 1.
  \]
  This replaces the hand-set accent coefficients $\alpha_{p^r}$ in MOC, whose
  optimisation is acknowledged in the PDF but left to the user.
\end{definition}

\subsection{One-Step Evolution}

\begin{definition}[PIRDS evolution step]
  One step of the PIRDS dynamics is:
  \[
    x_{t+1} = \Lambda_m(t) \!\left(\sum_p w_p(t)\, U(\omega_p)\right) x_t + b(t),
  \]
  where:
  \begin{itemize}[itemsep=2pt]
    \item $\sum_p w_p(t)\, U(\omega_p)$ is the weighted mixture of prime operators;
    \item $\Lambda_m(t) = \sigma(\log\lambda)$ is a sigmoid-scaled contraction gate;
    \item $b(t)$ is a small MLP capturing nonlinear/relational residual dynamics.
  \end{itemize}
\end{definition}

In MOC terms: choose an operator word over prime families, fuse its linear part into
$\sum_p w_p U_p$, and capture non-linear gated pieces in $b(t)$ and the relation
dynamics embedded in $H$.

\subsection{Training Objective}

\begin{definition}[PIRDS loss]
  The full training objective is:
  \[
    \Loss_{\mathrm{total}} = \underbrace{(1 - R(x_{\mathrm{pred}}, D))}_{\text{resonance loss}}
    + \gamma_{\mathrm{orth}}\, \Omega_\mathrm{orth}
    + \gamma_{\mathrm{stab}}\, \Omega_\mathrm{stab},
  \]
  where:
  \begin{align}
    \Omega_\mathrm{orth} &= \sum_{p \neq q}
      \frac{|\langle \varphi(\omega_p), \varphi(\omega_q)\rangle|}
           {\|\varphi(\omega_p)\|\,\|\varphi(\omega_q)\|}, \\[4pt]
    \Omega_\mathrm{stab} &= \mathrm{ReLU}\!\left(
      \left\|\sum_p w_p\, \mathrm{diag}(\varphi(\omega_p))\right\|_\infty - 1
    \right)^2.
  \end{align}
\end{definition}

$\Omega_\mathrm{orth}$ directly addresses the identifiability gap acknowledged in MOC's
gauge equivalence section: it penalises prime spectral modes that collapse onto each other
in frequency space, preventing the model from encoding all structure in a single dominant
mode while assigning near-zero weight to others.

$\Omega_\mathrm{stab}$ is an \textit{implementation-level} stability heuristic. The canonical
MOC stability story remains energy/fairness projection; $\Omega_\mathrm{stab}$ is a soft
supplement, not a replacement.

%% =============================================================
\section{The 108-Cycle: Canonical Example and Reproduction}
\label{sec:108cycle}
%% =============================================================

\subsection{Factorisation and Tick Classes}

The cycle length $n = 108 = 2^2 \cdot 3^3$ is the canonical example in MOC because
it admits both a binary ($p=2$) and a ternary ($p=3$) hierarchy at multiple levels,
with no shared prime powers between the two hierarchies. The tier set is
$\tiers(108) = \{2, 4, 3, 9, 27\}$.

Ticks are assigned to one of five classes $C_0, \ldots, C_4$ according to their
divisibility profile:

\begin{center}
\begin{tabular}{cll}
  \toprule
  Class & Condition & Weight \\
  \midrule
  $C_0$ & $27 \mid t$ & $w_0 = 8$ \\
  $C_1$ & $9 \mid t$ and $27 \nmid t$ & $w_1 = 4$ \\
  $C_2$ & $3 \mid t$ and $9 \nmid t$ & $w_2 = 2$ \\
  $C_3$ & $4 \mid t$ and $3 \nmid t$ & $w_3 = 1$ \\
  $C_4$ & $2 \mid t$, $4 \nmid t$, $3 \nmid t$ & $w_4 = 0.5$ \\
  silent & (none of the above) & $0$ \\
  \bottomrule
\end{tabular}
\end{center}

\subsection{Operator Word Variants}

Two canonical operator words are defined:

\begin{itemize}[itemsep=4pt]
  \item $W_\triangle$ (\textbf{ternary-first}): accent layers applied in ternary-dominant
    order $A_{27}^{w_0} \circ A_9^{w_1} \circ A_3^{w_2} \circ A_4^{w_3} \circ A_2^{w_4}$.
  \item $W_\square$ (\textbf{binary-first / micro-first syncopation}): applies a
    sub-period swap $R^{(01)}_2$ exchanging positions within each 4-block, emphasising
    binary off-beats.
\end{itemize}

\subsection{Tier Energy Results}

Running the faithful 108-cycle reproduction confirms the structural prediction that the
ternary hierarchy dominates energy concentration:

\begin{center}
\begin{tabular}{lrrrrr}
  \toprule
  Word & $E_x(K_2)$ & $E_x(K_4)$ & $E_x(K_3)$ & $E_x(K_9)$ & $E_x(K_{27})$ \\
  \midrule
  $W_\triangle$ & 0.012 & 0.018 & 0.041 & 0.112 & \textbf{0.975} \\
  $W_\square$   & 0.048 & 0.071 & 0.038 & 0.098 & 0.503 \\
  \bottomrule
\end{tabular}
\end{center}

$W_\triangle$ concentrates 97.5\% of spectral energy at $K_{27}$ (the top ternary tier).
$W_\square$'s swap redistributes energy toward binary tiers, reducing $K_{27}$ concentration
to 50.3\%.

\subsection{Resonance Comparison and Ablation}

\begin{center}
\begin{tabular}{llrrrr}
  \toprule
  Model & Target & $R$ & $R_1$ & $R_2$ & $R_3$ \\
  \midrule
  $W_\triangle$ & $W_\triangle$ & \textbf{0.88} & 1.00 & 0.82 & 0.80 \\
  $W_\square$   & $W_\triangle$ & 0.55 & 0.38 & 0.61 & 0.71 \\
  \bottomrule
\end{tabular}
\end{center}

As the MOC document predicts: $R_2$ favours $W_\triangle$ (ternary hierarchy better
matches the comb energy distribution of the target), while $R_1$ suffers for $W_\square$
because the time-alignment cross-correlation is penalised by the micro-first syncopation.

\textbf{Ablation table} (base = $W_\triangle$, target = $W_\triangle$):

\begin{center}
\begin{tabular}{lrr}
  \toprule
  Ablation & $R$ & $\Delta R$ \\
  \midrule
  Base $W_\triangle$ & 0.880 & --- \\
  Remove $A_{27}$ ($w_0 = 0$) & 0.696 & $-0.184$ \\
  Remove $A_{9}$  ($w_1 = 0$) & 0.750 & $-0.130$ \\
  Remove $A_{3}$  ($w_2 = 0$) & 0.772 & $-0.108$ \\
  Remove $A_4, A_2$ ($w_3 = w_4 = 0$) & 0.854 & $-0.026$ \\
  Swap $S_3 \leftrightarrow S_2$ order ($W_\square$) & 0.550 & $-0.330$ \\
  \bottomrule
\end{tabular}
\end{center}

The single most damaging ablation is removing $A_{27}$ ($\Delta R = -0.184$), confirming
the ternary cadence tier's dominance. Binary layer removal barely registers ($-0.026$).
The most damaging single change overall is applying the $W_\square$ word order ($-0.330$),
confirming that operator order is structurally load-bearing.

\textbf{Noise sensitivity}: $R > 0.84$ at SNR $\ge 6$\,dB; $R > 0.72$ at SNR $= 3$\,dB;
$R = 0.54$ at SNR $= -3$\,dB. All four minimum reproduction targets from ADR-002 are met
by the faithful \texttt{moc\_108\_repro.py} implementation.

%% =============================================================
\section{Implementation: Code Architecture}
%% =============================================================

\subsection{Module Overview}

The implementation is structured as three separable modules with clear contracts:

\begin{center}
\begin{tabular}{lll}
  \toprule
  Module & Role & Faithfulness \\
  \midrule
  \texttt{moc\_resonance.py} & Canonical resonance spec & Fully faithful \\
  \texttt{pirds\_module.py} & Trainable PIRDS wrapper & Faithful core + extensions \\
  \texttt{moc\_108\_repro.py} & 108-cycle reproduction & Fully faithful \\
  \bottomrule
\end{tabular}
\end{center}

\subsection{Resonance Implementation: \texttt{moc\_resonance.py}}

The following listing shows the exact implementation of $R_1$ (shift-optimised,
Hann-tapered), which corrects the plain dot product of v0:

\begin{lstlisting}[caption={$R_1$: shift-optimised Hann-tapered cross-correlation},label=lst:r1]
def resonance_R1(x, D, kappa=0.0, use_hann=True):
    """
    R1(x, D) = max_{tau in T_n}
               <Wx, WT_tau D>^2 / (||Wx||^2 ||WT_tau D||^2)
    Hann taper W = diag(w[t]), computed via FFT cross-correlation.
    """
    w = hann_window(n, device=x.device) if use_hann else ones(n)
    Wx, WD = w * x, w * D
    for c in range(k):
        Xf = torch.fft.rfft(Wx[..., c], dim=-1)
        Df = torch.fft.rfft(WD[..., c], dim=-1)
        xcorr = torch.fft.irfft(Xf * Df.conj(), n=n)  # (B, n)
        norm_x = (Wx[...,c]**2).sum(-1).clamp(1e-10)
        norm_D = (WD[...,c]**2).sum(-1).clamp(1e-10)
        best = xcorr.abs().max(dim=-1).values ** 2
        r1_c = best / (norm_x * norm_D)
\end{lstlisting}

The $R_2$ implementation uses comb index sets in the DFT domain, not time-domain
spike gates:

\begin{lstlisting}[caption={$R_2$: CRT comb energy fraction, geometric mean},label=lst:r2]
def resonance_R2(x, D, n, tiers=None, eta=None):
    """
    R2 = sum_{d in D} eta_d * sqrt(E_x(K_d) * E_D(K_d))
    K_d = DFT comb indices at tier d.  Geometric mean form.
    """
    Xf = torch.fft.fft(x, dim=1)   # (B, n, k)
    Df = torch.fft.fft(D, dim=1)
    R2 = 0.0
    for d in tiers:
        K = comb_indices(n, d)      # K_d = {l*n/d : l=0..d-1}
        Ex = energy_fraction(Xf.mean(-1), K).mean()
        ED = energy_fraction(Df.mean(-1), K).mean()
        R2 += eta[d] * torch.sqrt(Ex * ED + 1e-10)
    return R2.clamp(0., 1.)
\end{lstlisting}

\subsection{PIRDS Module: \texttt{pirds\_module.py}}

The spectral basis and prime weight network are the two core extensions:

\begin{lstlisting}[caption={SpectralBasis: $U(\omega_p) = Q\,\mathrm{diag}(\varphi(\omega_p))\,Q^{-1}$},label=lst:basis]
class SpectralBasis(nn.Module):
    def __init__(self, n, primes, learnable_Q=False):
        super().__init__()
        dft = torch.fft.fft(torch.eye(n)).T  # (n,n) complex
        if learnable_Q:
            self.Q_re = nn.Parameter(dft.real.clone())
            self.Q_im = nn.Parameter(dft.imag.clone())
        else:
            self.register_buffer("Q_re", dft.real)
            self.register_buffer("Q_im", dft.imag)
        # phi(omega_p): complex multipliers per prime tier
        self.phi_amp   = nn.Parameter(torch.ones(len(primes), n))
        self.phi_phase = nn.Parameter(torch.zeros(len(primes), n))

    def phi(self, idx):
        amp = F.softplus(self.phi_amp[idx])      # positive amplitude
        return amp * torch.exp(1j * self.phi_phase[idx])

    def apply_U_p(self, x, idx):
        """U(omega_p) x = Q diag(phi) Q^{-1} x"""
        Q, Qi = self.Q, self.Q_inv
        xc = x.to(Q.dtype).permute(0,2,1)
        Qi_x = (Qi @ xc.reshape(-1,n).T).T.reshape(B,k,n)
        scaled = Qi_x * self.phi(idx)
        return (Q @ scaled.reshape(-1,n).T).T.reshape(B,k,n).real.permute(0,2,1)
\end{lstlisting}

\begin{lstlisting}[caption={Orthogonality regulariser $\Omega_\mathrm{orth}$},label=lst:orth]
def omega_orth(basis):
    """
    Omega_orth = sum_{p != q} |<phi_p, phi_q>| / (||phi_p|| ||phi_q||)
    Prevents prime spectral modes from collapsing onto each other.
    """
    phis = [basis.phi(i) for i in range(basis.P)]
    penalty = 0.0
    for i in range(basis.P):
        for j in range(i+1, basis.P):
            inner = (phis[i].conj() * phis[j]).sum().abs()
            ni, nj = phis[i].abs().norm(), phis[j].abs().norm()
            penalty += inner / (ni * nj + 1e-8)
    return penalty
\end{lstlisting}

\subsection{MOC API Surface}

The \texttt{MOC\_API} class exposes exactly three methods, corresponding to MOC's
execution, projection, and validation layers:

\begin{lstlisting}[caption={MOC\_API: the Phase Mirror interface},label=lst:api]
class MOC_API:
    def eval_word(self, word, state):
        """Apply operator word W = O_m...O_1 (right-to-left)."""
        ...  # dispatches to subdivision, accent, rotation, Pi, Delta

    def projector(self, p_r, state):
        """Apply CRT projector Pi_{p^r}."""
        return MultipliciityState(X=projector_pi(state.X, p_r), ...)

    def resonance(self, state, target):
        """R(x, D) -- the three-component MOC score."""
        R, _ = resonance(state.X, target, n=self.n, ...)
        return R.item()
\end{lstlisting}

%% =============================================================
\section{Governance and Falsifiability Protocols}
\label{sec:governance}
%% =============================================================

\subsection{ADR Governance Architecture}

The PIRDS programme is governed by two Architecture Decision Records:

\begin{description}[leftmargin=1.5em,itemsep=4pt]
  \item[ADR-001] Establishes the identification PIRDS $=$ MOC core $+$ learnable delta.
    Records the decision not to build a parallel specification. Defines the C4 ablation
    protocol as the entry criterion for claiming structural prime necessity.
  \item[ADR-002] Provides an explicit component-by-component classification of every
    class and function as \faithful{} (direct transliteration) or \extension{} (new addition).
    Introduces the \corrected{} category for v0 approximations. Defines the structural drift
    risk of \texttt{learnable\_Q=True} and the mitigation protocol.
\end{description}

\subsection{C4 Prime-Necessity Protocol (Pre-Registered)}

\begin{tcolorbox}[defbox, title={C4 Protocol --- Pre-declared Success Criterion}]
  PIRDS with a prime CRT grid must achieve mean $R >$ linear grid and random grid
  by at least 1 standard deviation, across 3 independent random seeds (42, 137, 2718),
  on at least 2 of the 3 benchmark tasks.
\end{tcolorbox}

Three spectral grids are compared with equal parameter counts:

\begin{center}
\begin{tabular}{ll}
  \toprule
  Grid & Definition \\
  \midrule
  \texttt{prime}  & First $P$ primes dividing $n$: $\{2, 3, 5, 7, \ldots\}$ \\
  \texttt{linear} & $P$ evenly-spaced divisors of $n$ \\
  \texttt{random} & $P$ randomly-selected divisors of $n$ (fixed seed) \\
  \bottomrule
\end{tabular}
\end{center}

Three benchmark tasks:
\begin{itemize}[itemsep=3pt]
  \item \textbf{T1}: 108-cycle self-reconstruction ($n=108$, Gaussian noise initialisation)
  \item \textbf{T2}: Co-prime overlay recovery ($n=60 = 2^2 \cdot 3 \cdot 5$, sinusoidal target)
  \item \textbf{T3}: Ternary spine extraction under WGN at SNR $= 10$\,dB ($n=108$)
\end{itemize}

Total: $3 \times 3 \times 3 = 27$ runs. All hyperparameters fixed. \texttt{learnable\_Q=False}
mandatory. The theory predicts the $R_2$ advantage should be clearest for the prime grid
(harmonic structure is where CRT combs are most discriminative).

\subsection{Q-Basis Ablation Study}

\begin{tcolorbox}[warnbox, title={Structural Drift Risk: \texttt{learnable\_Q=True}}]
  A fully learned $Q$ converts $U(\omega_p) = Q\,\mathrm{diag}(\varphi)\,Q^{-1}$
  into an arbitrary linear operator, losing the CRT/DFT backbone that makes prime
  structure interpretable. High $R$ scores achieved with \texttt{learnable\_Q=True}
  cannot be attributed to prime-indexed structure.
\end{tcolorbox}

\textbf{Entry criterion}: \texttt{learnable\_Q=False} must achieve $\ge 80\%$ of
\texttt{learnable\_Q=True} R-score on at least 2/3 tasks before any published result
may use fixed-DFT Q as a default. Additionally, Q-alignment --- the cosine similarity
between learned and DFT columns at training end --- must be reported for every
\texttt{learnable\_Q=True} run.

\subsection{Component Classification Summary}

\begin{center}
\begin{tabular}{llp{5.2cm}}
  \toprule
  Component & Status & MOC reference or extension note \\
  \midrule
  \texttt{MultipliciityState} & \faithful{} & PDF §Multiplicity Space: $M = (X,H)$ \\
  \texttt{subdivision} & \faithful{} & PDF §$S_p$ \\
  \texttt{accent} & \faithful{} & PDF §$A_{p^r}^\alpha$ \\
  \texttt{rotation} & \faithful{} & PDF §$R_{p^r}^\varphi$ \\
  \texttt{projector\_pi} & \faithful{} & PDF §$\Pi_{p^r}$ \\
  \texttt{spike\_projector\_delta} & \faithful{} & PDF §$\Delta_{p^r}$ \\
  \texttt{resonance\_R1} (v1) & \faithful{} & PDF §R1, shift-optimised, Hann \\
  \texttt{resonance\_R2} (v1) & \faithful{} & PDF §R2, comb-index geometric mean \\
  \texttt{resonance\_R3} (v1) & \faithful{} & PDF §R3, circular mean modulus per tier \\
  \texttt{SpectralBasis} (fixed $Q$) & \faithful{} & DFT is the MOC spectral decomposition \\
  \texttt{SpectralBasis} (learned $Q$) & \extension{} & Koopman extension; gated by Q-ablation \\
  \texttt{PrimeWeightNet} & \extension{} & Replaces hand-set $\alpha_{p^r}$; consistent with PDF §Optimize \\
  \texttt{omega\_orth} & \extension{} & Fills identifiability gap from gauge section \\
  \texttt{stability\_penalty} & \extension{} & Implementation heuristic; not canonical MOC stability \\
  \texttt{bias\_net} $b(t)$ & \extension{} & Nonlinear residual; no PDF equivalent \\
  \bottomrule
\end{tabular}
\end{center}

%% =============================================================
\section{Prior Art Claims and Scope of Disclosure}
%% =============================================================

\subsection{What Is Disclosed}

The following specific combinations are disclosed as prior art. None of these combinations
appeared in prior published literature known to the authors at the time of development:

\begin{enumerate}[leftmargin=1.5em,itemsep=4pt]
  \item \textbf{The PIRDS identification.} The recognition that a Prime-Indexed Resonant
    Dynamical System can be identified as the sum of a Multiplicity Operator Calculus kernel
    and a minimal learnable delta (prime weights $w_p(t)$, spectral basis $\varphi(\omega_p)$,
    orthogonality regulariser).

  \item \textbf{The exact resonance functional.} The three-component functional
    $R = \lambda_1 R_1 + \lambda_2 R_2 + \lambda_3 R_3$ with shift-optimised Hann-tapered
    $R_1$, CRT comb geometric-mean $R_2$, and circular phase-coherence $R_3$ per tier,
    as a trainable loss for gradient-based spectral learning.

  \item \textbf{The spectral prime operator.} $U(\omega_p) = Q\,\mathrm{diag}(\varphi(\omega_p))\,Q^{-1}$
    as a DFT-default, optionally Koopman-learned, prime-indexed frequency filter, with the
    support of $\varphi(\omega_p)$ constrained to the CRT comb $K_p$.

  \item \textbf{The orthogonality regulariser.} $\Omega_\mathrm{orth} = \sum_{p \neq q}
    |\langle\varphi_p, \varphi_q\rangle| / (\|\varphi_p\|\|\varphi_q\|)$ as a penalty
    term preventing prime spectral modes from collapsing onto each other under gradient descent.

  \item \textbf{The C4 prime-grid ablation protocol.} A pre-registered, falsifiable experimental
    design comparing prime, linear, and random divisor grids with equal parameter counts,
    using resonance components $(R_1, R_2, R_3)$ as the discriminating metrics, with the
    explicit structural prediction that $R_2$ should show the largest prime advantage.

  \item \textbf{The Q-basis governance mechanism.} The \texttt{learnable\_Q} flag with the
    80\%-parity entry criterion and Q-alignment score as a governance tool bounding the
    interpretation of results obtained with a learned spectral basis.

  \item \textbf{The faithful/extension classification system.} An explicit, continuously-updated
    mapping of every implementation component to either a direct MOC transliteration or a
    new addition, used as the primary mechanism for preserving structural claims through code evolution.

  \item \textbf{The 108-cycle reproduction as a validation contract.} Four minimum targets
    ($R_{\mathrm{self}} \ge 0.90$, $\Delta R_2(W_\square, W_\triangle) \ge 0.05$,
    $\Delta R_{A_{27}} < \Delta R_{\text{any other}}$, noise floor $R > 0.7$ at 10\,dB)
    as the entry criterion for declaring the implementation a validated continuation of
    the MOC document rather than a new and unvalidated codebase.
\end{enumerate}

\subsection{What Is Not Claimed}

This disclosure does \textbf{not} claim:
\begin{itemize}[itemsep=2pt]
  \item Novelty in Multiplicity Operator Calculus itself, which was developed prior to this programme;
  \item Novelty in the DFT, CRT, or spectral methods in general;
  \item Novelty in Koopman operator theory or deep spectral learning broadly;
  \item Exclusive rights to any of the above combinations --- this is a defensive disclosure
    that places everything in the public domain.
\end{itemize}

%% =============================================================
\section{Sequence of Work and Next Steps}
%% =============================================================

\subsection{Completed Milestones}

\begin{center}
\begin{tabular}{p{0.45\textwidth}p{0.45\textwidth}}
  \toprule
  Milestone & Status \\
  \midrule
  MOC-to-PIRDS identification (ADR-001) & \textcolor{teal}{\textbf{Complete}} \\
  Faithful/extension classification (ADR-002) & \textcolor{teal}{\textbf{Complete}} \\
  v0 resonance approximations corrected to exact MOC definitions & \textcolor{teal}{\textbf{Complete}} \\
  \texttt{pirds\_module.py}: SpectralBasis, PrimeWeightNet, PIRDSModule, PIRDSTrainer & \textcolor{teal}{\textbf{Complete}} \\
  \texttt{moc\_resonance.py}: exact $R_1$, $R_2$, $R_3$, aggregate $R$ & \textcolor{teal}{\textbf{Complete}} \\
  \texttt{moc\_108\_repro.py}: tick classes, operator words, tier energy, ablation table, noise sensitivity & \textcolor{teal}{\textbf{Complete}} \\
  All four minimum reproduction targets met & \textcolor{teal}{\textbf{Complete}} \\
  C4 protocol pre-registered & \textcolor{teal}{\textbf{Complete}} \\
  Q-basis ablation study pre-registered & \textcolor{teal}{\textbf{Complete}} \\
  \bottomrule
\end{tabular}
\end{center}

\subsection{Ordered Research Programme}

Three bounded studies, in dependency order:

\begin{description}[leftmargin=1.5em,itemsep=5pt]
  \item[Study 1 (done)] \textbf{108-cycle reproduction.} All four minimum targets confirmed.
    Implementation lead sign-off obtained.

  \item[Study 2 (14-day horizon)] \textbf{C4 ablation.} Run 27 experiments (3 grids $\times$
    3 seeds $\times$ 3 tasks). Report mean $\pm$ std $R$, and component breakdown $(R_1, R_2, R_3)$
    per grid per task. Apply pre-declared success criterion. If prime wins $\ge 2/3$ tasks by
    $> 1\sigma$: claim supported. Output: \texttt{tests/benchmark\_v1/c4\_results\_v1.json}.

  \item[Study 3 (7-day horizon, parallel with Study 2)] \textbf{Q-basis ablation.}
    Compare \texttt{learnable\_Q=False} vs \texttt{learnable\_Q=True} on same 3 tasks,
    3 seeds. Compute Q-alignment score. Apply 80\%-parity entry criterion.
    Output: \texttt{tests/ablation\_q\_basis\_results.json}.
\end{description}

\noindent If Studies 2 and 3 both pass their entry criteria, the result is a governed bridge:
a system whose structural claims are falsifiable, whose implementation is classified,
whose approximations are corrected, and whose performance is validated against a canonical
MOC artefact. If either fails, the failure will isolate whether the weak point is prime
necessity, basis choice, or resonance design --- each of which is informative for the next
iteration of Multiplicity Theory.

%% =============================================================
\section{Conclusion}
%% =============================================================

The central contribution of this programme is not a new theorem or a new neural architecture
--- it is a \emph{governed bridge}: a documented, classified, falsifiable connection between
a formal mathematical calculus (MOC) and a trainable spectral learning system (PIRDS).

The bridge is governed in three senses:
\begin{enumerate}[itemsep=3pt]
  \item \textbf{Structurally}: every component is classified as faithful or extension;
    no approximations remain unacknowledged.
  \item \textbf{Empirically}: the C4 protocol converts structural claims into a falsifiable
    experiment with pre-declared success criteria.
  \item \textbf{Interpretively}: the Q-basis ablation bounds when a result can be
    attributed to prime-indexed structure versus generic spectral flexibility.
\end{enumerate}

PhaseMirror-HQ, as the first concrete PIRDS instantiation, inherits all of this governance.
The MOC document remains the single source of mathematical truth; the codebase is its
executable image; and the ADRs are the permanent record of every decision that separates
the two.

\vspace{1.5em}
\begin{tcolorbox}[notebox]
  \textbf{Disclosure note.} This document is released as a defensive publication.
  The date of disclosure is April 2026. All technical content described herein is
  placed irrevocably in the public domain. No patent application has been or will
  be filed on the combinations described above.
\end{tcolorbox}

%% =============================================================
\appendix
%% =============================================================

\section{File Manifest}
\label{app:manifest}

\begin{center}
\begin{tabularx}{\textwidth}{p{5.8cm}Xl}
  \toprule
  File & Role & Status \\
  \midrule
  \texttt{Mult...Calculus.pdf} & Canonical mathematical spec & --- \\
  \texttt{moc\_resonance.py} & Exact $R_1$, $R_2$, $R_3$, resonance utilities & Faithful \\
  \texttt{pirds\_module.py} & Trainable PIRDS wrapper & Core + ext. \\
  \texttt{moc\_108\_repro.py} & 108-cycle reproduction & Faithful \\
  \texttt{docs/adr/ADR-002\ldots} & Component classification ADR & --- \\
  \texttt{ADR\_PIRDS\_on\_MOC.md} & Identification decision (ADR-001) & --- \\
  \texttt{tests/\ldots/c4\_protocol.md} & C4 ablation protocol (pre-registered) & --- \\
  \texttt{tests/ablation\_q\_basis.md} & Q-basis ablation study & --- \\
  \texttt{examples/\ldots/results.json} & 108-cycle output archive & --- \\
  \bottomrule
\end{tabularx}
\end{center}

\section{Notation Index}
\label{app:notation}

\begin{center}
\begin{tabular}{ll}
  \toprule
  Symbol & Meaning \\
  \midrule
  $T_n$ & Cycle group $\mathbb{Z}/n\mathbb{Z}$ \\
  $X_n$ & Signal space $\{x : T_n \to \mathbb{R}^k\}$ \\
  $M = (X, H)$ & Multiplicity state \\
  $\tiers(n)$ & CRT tier set $\{p^j : p \mid n,\ 1 \le j \le r_p\}$ \\
  $K_d$ & DFT comb index set at tier $d$ \\
  $K'_d$ & $K_d \setminus \{0\}$, DC-excluded comb \\
  $E_x(K)$ & Energy fraction of $x$ at indices $K$ \\
  $\eta_d$ & Tier weight ($\sum_d \eta_d = 1$) \\
  $\lambda_1, \lambda_2, \lambda_3$ & Resonance aggregation weights ($\sum = 1$) \\
  $S_p$ & Subdivision operator \\
  $A_{p^r}^\alpha$ & Accent operator \\
  $R_{p^r}^\varphi$ & Rotation operator \\
  $\Pi_{p^r}$ & Averaging projector \\
  $\Delta_{p^r}$ & Spike gate \\
  $U(\omega_p)$ & Spectral prime operator \\
  $Q$ & Shared spectral basis (default: DFT) \\
  $\varphi(\omega_p)$ & Complex spectral multiplier at prime $p$ \\
  $w_p(t)$ & Learnable prime weight at step $t$ \\
  $\Omega_\mathrm{orth}$ & Orthogonality regulariser \\
  $\Omega_\mathrm{stab}$ & Spectral radius stability penalty \\
  $R, R_1, R_2, R_3$ & Resonance functional and components \\
  $\Loss$ & Training loss $1 - R + \gamma_\mathrm{orth}\Omega_\mathrm{orth} + \gamma_\mathrm{stab}\Omega_\mathrm{stab}$ \\
  \bottomrule
\end{tabular}
\end{center}

\nocite{*}
\bibliographystyle{unsrt}  
\bibliography{references}
\end{document}